\subsection{Replies to comments received from Kate Scholberg - May 25, 2013}

\begin{enumerate}

%........................................................................................................
\vspace{5 mm}
\item
{\em
  I would like to see the F-C 1D contour created by the method Nakadaira-san suggests.
}\\
\\
We have performed a small 1D Feldman-Cousins production at the points of interest,
using the suggested method.

For the $\sin^2(2\theta_{23})$ limit, toys were generated with 
$|\Delta m^2_{32}| = 2.443 \times 10^-3$, and the input points and 90\% $\Delta \chi^2_{crit}$ values are:
\begin{itemize}
\item 0.945, 2.66;
\item 0.950, 2.56;
\item 0.955, 2.51.
\end{itemize}
These values are less than the constant-$\Delta \chi^2$ value of 2.71, and so would result in tighter limits.

For the $|\Delta m^2_{32}|$ lower limit, toys were generated with 
$\sin^2(2\theta_{23}) = 1.000$, and the input points and 68\% $\Delta \chi^2_{crit}$ values are:
\begin{itemize}
\item 2.28 $\times 10^{-3}$, 0.83;
\item 2.29 $\times 10^{-3}$, 0.89;
\item 2.30 $\times 10^{-3}$, 0.86;
\end{itemize}
These values are less than the constant-$\Delta \chi^2$ value of 1.00, and so would result in tighter limits.

For the $|\Delta m^2_{32}|$ upper limit, toys were generated with 
$\sin^2(2\theta_{23}) = 1.000$, and the input points and 68\% $\Delta \chi^2_{crit}$ values are:
\begin{itemize}
\item 2.60 $\times 10^{-3}$, 0.91;
\item 2.61 $\times 10^{-3}$, 0.85;
\item 2.62 $\times 10^{-3}$, 0.86;
\end{itemize}
These values are less than the constant-$\Delta \chi^2$ value of 1.00, and so would result in tighter limits.

%........................................................................................................
\vspace{5 mm}
\item
{\em
  One minor thing that occurred to me when rereading the memo:
  (I'm sure negligible, so do not worry about this for this iteration-- 
  there are much bigger effects not taken into account that we are neglecting-- 
  but perhaps consider for the future).
  Don't you also need an MC template for P(nue$\rightarrow$numu), for beam nue?  This is very tiny--by back-of-the-envelope,
  effect for r1+2+3 would be: for $\sim$1.5 beam nues (according to nue appearance bg estimate), and $\sim$10\% 
  osc effect at max, you'd get 0.1 oscillated numu in the sample (basically, a part-per-mil effect).
  But maybe should include for completeness.
}\\
\\
%We will not include a $\nu_{e} \rightarrow \nu_{\mu}$ template in the current Run 1+2+3+4a analysis, but 
We will consider including the additional templates in the final Run 1-4 analysis.  

\end{enumerate}
